\documentclass[11pt,a4paper]{report}
\usepackage{amsmath}
\usepackage{amssymb}
\usepackage{graphicx}
\usepackage{fontspec}
\usepackage{babel}
\usepackage[margin=3cm, top=3cm, bottom=3cm]{geometry}
\usepackage{enumitem}
\usepackage{booktabs}
\usepackage{xcolor}
\usepackage{hyperref}
\usepackage{setspace}
\usepackage{fancyhdr}

%% ─── Design system tokens (@docs/design/design.md) ──────────
\definecolor{primary}{HTML}{cc785c}
\definecolor{coral}{HTML}{cc785c}
\definecolor{primary-active}{HTML}{a9583e}
\definecolor{primary-disabled}{HTML}{e6dfd8}
\definecolor{ink}{HTML}{141413}
\definecolor{body-strong}{HTML}{252523}
\definecolor{body}{HTML}{3d3d3a}
\definecolor{muted}{HTML}{6c6a64}
\definecolor{muted-soft}{HTML}{8e8b82}
\definecolor{canvas}{HTML}{faf9f5}
\definecolor{surface-soft}{HTML}{f5f0e8}
\definecolor{surface-card}{HTML}{efe9de}
\definecolor{surface-cream-strong}{HTML}{e8e0d2}
\definecolor{surface-dark}{HTML}{181715}
\definecolor{surface-dark-elevated}{HTML}{252320}
\definecolor{surface-dark-soft}{HTML}{1f1e1b}
\definecolor{hairline}{HTML}{e6dfd8}
\definecolor{on-dark}{HTML}{faf9f5}
\definecolor{on-dark-soft}{HTML}{a09d96}
\definecolor{accent-teal}{HTML}{5db8a6}
\definecolor{accent-amber}{HTML}{e8a55a}
\definecolor{success}{HTML}{5db872}
\definecolor{warning}{HTML}{d4a017}
\definecolor{error}{HTML}{c64545}

\hypersetup{colorlinks=true, linkcolor=coral, urlcolor=coral}

\babelprovide[import, onchar=ids fonts]{thai}
\babelprovide[import, onchar=ids fonts]{english}

\babelfont{rm}{Noto Sans}
\babelfont[thai]{rm}{Noto Sans Thai}
\babelfont{tt}{Noto Sans Mono}
\babelfont[thai]{tt}{Noto Sans Thai}

\directlua{
  local glyph_id = node.id("glyph")
  local penalty_id = node.id("penalty")

  local function is_thai(c)
    return c >= 0x0E01 and c <= 0x0E5B
  end

  local function is_thai_combining(c)
    return (c >= 0x0E31 and c <= 0x0E3A) or (c >= 0x0E47 and c <= 0x0E4E)
  end

  luatexbase.add_to_callback("pre_linebreak_filter", function(head)
    for n in node.traverse_id(glyph_id, head) do
      local prev = n.prev
      if prev and prev.id == glyph_id and is_thai(prev.char)
         and is_thai(n.char) and not is_thai_combining(n.char) then
        local p = node.new(penalty_id)
        p.penalty = 0
        node.insert_before(head, n, p)
      end
    end
    return head
  end, "thai_break")
}

\emergencystretch=1em
\onehalfspacing

\pagestyle{fancy}
\fancyhf{}
\fancyhead[L]{\color{muted}\small บทที่ X — ชื่อบท}
\fancyhead[R]{\color{muted}\small\thepage}
\renewcommand{\headrulewidth}{0.3pt}

\begin{document}

%% ─── Chapter heading ──────────────────────────────────────────
\chapter{ทบทวนวรรณกรรม}

เนื้อหานำของบทนี้ อธิบายว่าบทนี้ครอบคลุมอะไรบ้าง
และเชื่อมโยงกับบทอื่นอย่างไร

%% ─── Section 1 ────────────────────────────────────────────────
\section{แนวคิดพื้นฐาน}

\subsection{แนวคิดที่ 1}

อธิบายแนวคิดพร้อมอ้างอิงงานวิจัยที่เกี่ยวข้อง

\subsection{แนวคิดที่ 2}

อธิบายแนวคิดพร้อมอ้างอิงงานวิจัยที่เกี่ยวข้อง

%% ─── Section 2 ────────────────────────────────────────────────
\section{งานวิจัยที่เกี่ยวข้อง}

\subsection{กลุ่มงานวิจัยที่ 1: ชื่อกลุ่ม}

สรุปงานวิจัยในกลุ่มนี้ จุดเด่น ข้อจำกัด
และความสัมพันธ์กับงานวิจัยของเรา

\subsection{กลุ่มงานวิจัยที่ 2: ชื่อกลุ่ม}

สรุปงานวิจัยในกลุ่มนี้ จุดเด่น ข้อจำกัด
และความสัมพันธ์กับงานวิจัยของเรา

%% ─── Section 3 ────────────────────────────────────────────────
\section{ช่องว่างของงานวิจัย}

\begin{itemize}[nosep, leftmargin=2em]
  \item ช่องว่างที่ 1: อธิบาย
  \item ช่องว่างที่ 2: อธิบาย
  \item ช่องว่างที่ 3: อธิบาย
\end{itemize}

\vspace{12pt}
งานวิจัยนี้เติมเต็มช่องว่างดังกล่าวโดย\ldots

%% ─── Section 4 ────────────────────────────────────────────────
\section{กรอบแนวคิดของงานวิจัย}

อธิบาย conceptual framework ที่ใช้ในงานวิจัยนี้
พร้อมแผนภาพหรือตารางถ้ามี

%% ─── Summary ──────────────────────────────────────────────────
\section{สรุปบท}

สรุปประเด็นสำคัญของบทนี้และนำไปสู่บทถัดไป

\end{document}
